\chapter{諸定理}
    \section{偶関数と奇関数への分解}
        \begin{shadebox}
            連続な1変数関数 $f$ は連続な偶関数と奇関数に 1 通りに分解できる
        \end{shadebox}
        \begin{proof}
            \quad\par
            (分解可能性)\\
            $\fe(x) \coloneq (f(x)+f(-x))/2,\fo(x) \coloneq (f(x)-f(-x))/2$ とすれば確かに $f(x)=\fe(x)+\fo(x)$ であり $\fe,\fo$ はそれぞれ連続な偶関数，奇関数になっている。\\
            \newline
            (一意性)
            \par
            次に，分解が 1 通りであることを示す。
            偶関数 $\fea,\feb$ と奇関数 $\foa,\fob$ が存在して $f(x) = \fea(x) + \foa(x) = \feb(x) + \fob(x)$ と表せるとする。
            このとき $\fea(x) - \feb(x) = \fob(x) - \foa(x)$ が成り立つ。
            左辺は偶関数，右辺は奇関数なので，両辺は恒等的に 0 でなければならない（偶対称かつ奇対称な関数が恒等的に 0 であることは容易に示せる）。
        \end{proof}
    \section{\texorpdfstring{$\sup_{\vx\in X}(f(\vx)+g(\vx)) \leq \sup_{\vx\in X}f(\vx) + \sup_{\vx\in X}g(\vx),\inf_{\vx\in X}f(\vx) + \inf_{\vx\in X}g(\vx) \leq \inf_{\vx\in X}(f(\vx)+g(\vx))$}{sup\_\{x∈X\} f(x)+g(x) ≦ sum\_\{x∈X\}f(x) + sup\_\{x∈X\}g(x), inf\_\{x∈X\}f(x) + inf\_\{x∈X\}g(x) ≦ inf\_\{x∈X\} f(x)+g(x)}}
        \begin{shadebox}
            集合$X$上で定義された有界な関数$f,g$について次が成り立つ。
            \[ \sup_{\vx\in X}(f(\vx)+g(\vx)) \leq \sup_{\vx\in X}f(\vx) + \sup_{\vx\in X}g(\vx),\inf_{\vx\in X}f(\vx) + \inf_{\vx\in X}g(\vx) \leq \inf_{\vx\in X}(f(\vx)+g(\vx)) \]
        \end{shadebox}
        \begin{proof}
            \quad\par
            $\sup_{\vx\in X}(f(\vx)+g(\vx))$を与える$\vx\in X$の一つを$\vx_1$とすると$\sup_{\vx\in X}(f(\vx)+g(\vx)) = f(\vx_1)+g(\vx_1) \leq \sup_{\vx\in X}f(\vx) + \sup_{\vx\in X}g(\vx)$。
            また，$\inf_{\vx\in X}(f(\vx)+g(\vx))$を与える$\vx\in X$の一つを$\vx_2$とすると$\inf_{\vx\in X}(f(\vx)+g(\vx)) = f(\vx_2)+g(\vx_2) \geq \inf_{\vx\in X}f(\vx) + \inf_{\vx\in X}g(\vx)$
        \end{proof}
    \section{\texorpdfstring{$\sup_{\vx\in X}(-f(\vx)) = -\inf_{\vx\in X}f(\vx), \inf_{\vx\in X}(-f(\vx)) = -\sup_{\vx\in X}f(\vx)$}{sup\_\{x∈X\} -f(x) = -inf\_\{x∈X\}f(x), inf\_\{x∈X\} -f(x) = -sup\_\{x∈X\}f(x)}}
        \begin{shadebox}
            集合$X$上で定義された有界な関数$f,g$について次が成り立つ。
            \[ \sup_{\vx\in X}(-f(\vx)) = -\inf_{\vx\in X}f(\vx), \inf_{\vx\in X}(-f(\vx)) = -\sup_{\vx\in X}f(\vx) \]
        \end{shadebox}
        \begin{proof}
            \quad\par
            $^\forall\vy\in X,\inf_{\vx\in X}f(\vx) \leq f(\vy)$なので$^\forall\vy\in X,-\inf_{\vx\in X}f(\vx) \geq -f(\vy)$。
            任意の$\varepsilon>0$に対して$\inf_{\vx\in X}f(\vx) + \varepsilon$は$f$の下界ではないから，適当な$\vx_1\in X$が存在して$f(\vx_1) < \inf_{\vx\in X}f(\vx) + \varepsilon \quad \therefore\; -f(\vx_1) > -\inf_{\vx\in X}f(\vx) - \varepsilon$
            。すなわち$-\inf_{\vx\in X}f(\vx) - \varepsilon$は$-f$の上界ではない。
            よって上限の定義より$\sup_{\vx\in X}(-f(\vx)) = -\inf_{\vx\in X}f(\vx)$。2つ目の主張の証明も同様。
        \end{proof}
    \section{H\"older(ヘルダー)の不等式}
        \begin{shadebox}
            $p,q>1,\frac{1}{p}+\frac{1}{q}=1$のとき
            \[\sum_{i=1}^n|x_iy_i|\leq\norm{\vx}_p\norm{\vy}_q\]
        \end{shadebox}
        \begin{proof}
            \quad\par
            $\vx,\vy$の少なくとも1つが零ベクトルであるときは明らかに成り立つ。
            以下では両者とも零ベクトルでないとする。
            Youngの不等式より$a,b\geq 0$に対して$ab \leq a^p/p + b^q/q$が成り立つ。
            この式で$a=|x_i|/\norm{\vx}_p,b=|y_i|/\norm{\vy}_q$とすると
            \[\frac{|x_iy_i|}{\norm{\vx}_p\norm{\vy}_q}\leq\frac{|x_i|^p/\norm{\vx}_p^p}{p}+\frac{|y_i|^q/\norm{\vy}_q^q}{q}\]
            これを$i=1,\dots,n$まで足し合わせると
            \[\frac{1}{\norm{\vx}_p\norm{\vy}_q}\sum_{i=1}^n{|x_iy_i|} \leq \frac{\norm{\vx}_p^p/\norm{\vx}_p^p}{p}+\frac{\norm{\vy}_q^q/\norm{\vy}_q^q}{q} = 1/p+1/q = 1\]
            \[\therefore\;\sum_{i=1}^n|x_iy_i|\leq\norm{\vx}_p\norm{\vy}_q\]
        \end{proof}
    \section{Minkowskiの不等式(三角不等式の一般化)}
        \begin{shadebox}
            $1<p<\infty$のとき
            \[\norm{\vx+\vy}_p \leq \norm{\vx}_p+\norm{\vy}_p\]
        \end{shadebox}
        \begin{proof}
            \begin{equation}
                \norm{\vx+\vy}_p^p = \sum_{i=1}^n{|x_i+y_i|^p} = \sum_{i=1}^n{|x_i+y_i||x_i+y_i|^{p-1}} \leq \sum_{i=1}^n{|x_i||x_i+y_i|^{p-1}}+\sum_{i=1}^n{|y_i||x_i+y_i|^{p-1}}
            \end{equation}
            右辺第1項はH\"olderの不等式より上から
            \[\norm{\vx}_p\norm{\vv}_{\frac{p}{p-1}}\quad\mathrm{where}\quad\vv[i]\coloneq|x_i+y_i|^{p-1}\]
            で抑えられる。さらに
            \[\norm{\vx}_p\norm{\vv}_{\frac{p}{p-1}} = \norm{\vx}_p\bracks*{\sum_{i=1}^n{|x_i+y_i|^{(p-1)\frac{p}{p-1}}}}^{\frac{p-1}{p}} = \norm{\vx}_p\norm{\vx+\vy}_p^{p-1}\]
            同様に式(1)の右辺第2項も上から$\norm{\vy}_p\norm{\vx+\vy}_p^{p-1}$で抑えられるので
            \[\norm{\vx+\vy}_p^p \leq \norm{\vx}_p\norm{\vx+\vy}_p^{p-1} + \norm{\vy}_p\norm{\vx+\vy}_p^{p-1}\]
            両辺を$\norm{\vx+\vy}_p^{p-1}$で割ることでMinkowskiの不等式を得る。
        \end{proof}